\documentclass[a4paper, 14pt]{article}
\usepackage[margin=1.6cm]{geometry}
\usepackage[utf8]{inputenc}
\usepackage{minted}
\usepackage[russian]{babel}
\usepackage{amsmath}
\usepackage{graphicx}
\usepackage{changepage}
\usepackage{hyperref}
\usepackage{cases}
\usepackage{tikz-timing}[2017/12/20]
\usepackage{relsize}
\usepackage{booktabs}
\usepackage{gensymb}
\usepackage{multirow}
\usepackage{longtable}
\usetikzlibrary {arrows.meta}

\hypersetup{
	linkbordercolor = {1 1 1}
}

\usepackage{tikz-timing}[2009/05/15]
\usepackage{multicol}
\usepackage[T2A]{fontenc}
\usepackage{pgfplots}
%\usepackage[left=2.5cm, right=1.5cm, vmargin=2.5cm]{geometry}
\setlength\parindent{0pt} % Удалить отступы из параграфов.

\usepackage{listings}
\usepackage{caption}
\DeclareCaptionFont{white}{\color{white}} % Текст заголовка.
\DeclareCaptionFormat{listing}{\colorbox{gray}{\parbox{\textwidth}{#1#2#3}}}
\captionsetup[lstlisting]{format=listing,labelfont=white,textfont=white}
\renewcommand\labelenumi{\theenumi)}
\setlength\parindent{24pt}



\begin{document}
\lstset{
    language=java,                 % Выбор языка для подсветки (здесь это java).
    basicstyle=\small\sffamily,    % Размер и начертание шрифта для подсветки кода.
    numbers=left,                  % Где поставить нумерацию строк (слева\справа).
    numberstyle=\tiny,             % Размер шрифта для номеров строк.
    stepnumber=1,                  % Размер шага между двумя номерами строк.
    firstnumber=1,
    numberfirstline=true
    numbersep=5pt,                 % Как далеко отстоят номера строк от подсвечиваемого кода.
    backgroundcolor=\color{white}, % Цвет фона подсветки - используем \usepackage{color}.
    showspaces=false,              % Показывать или нет пробелы специальными отступами.
    showstringspaces=false,        % Показывать или нет пробелы в строках.
    showtabs=false,                % Показывать или нет табуляцию в строках.
    frame=single,                  % Рисовать рамку вокруг кода.
    tabsize=2,                     % Размер табуляции по умолчанию равен 2 пробелам.
    captionpos=t,                  % Позиция заголовка вверху [t] или внизу [b].
    breaklines=true,               % Автоматически переносить строки (да\нет).
    breakatwhitespace=false,       % Переносить строки только если есть пробел.
    escapeinside={\%*}{*)}         % Если нужно добавить комментарии в коде.
}

\begin{titlepage}
    \center

    ФЕДЕРАЛЬНОЕ ГОСУДАРСТВЕННОЕ АВТОНОМНОЕ ОБРАЗОВАТЕЛЬНОЕ УЧРЕЖДЕНИЕ ВЫСШЕГО ОБРАЗОВАНИЯ\linebreak
    «Санкт-Петербургский политехнический университет Петра Великого»
    \noindent\rule{500pt}{0.8pt} \\
    \textsc{\Large Институт компьютерных наук и кибербезопасности}\\
    \textsc{\large Высшая школа программной инженерии}\\[1.5cm]

    { \huge \bfseries ОТЧЕТ ПО ИНТЕГРАЦИОННОМУ ТЕСТИРОВАНИЮ	\\
    \Large \mdseries АГРЕГАТОР ЦИФРОВЫХ ФИНАНСОВЫХ АКТИВОВ <<ТЕССЕРАКТ>> \\
    \large по дисциплине <<Технологии разработки качественного программного обеспечения>>}\\
    \flushright{
        {\phantom{qwe}}\\[1.0cm]
    }

    \begin{figure}[H]
        \centering
        \includegraphics[width=10cm]{./resources/1.png}\\[2.0cm]
    \end{figure}

    \begin{multicols}{2}
        \begin{flushright} \large

            {Выполнили студенты группы: 5130904/00104:}\\
            {\phantom{qwe}}\\
            {\phantom{qwe}}\\
            {\phantom{qwe}}\\
            {\phantom{qwe}}\\

            {Преподаватель:\\}

        \end{flushright}
        \begin{flushright}

            {Почернин В. С.}\\
            {Шиляев В. С.}\\
            {Мурзаканов И. М.}\\
            {Разукрантов В. Е.}\\[0.5cm]


            Маслаков А. П.\\

        \end{flushright}
    \end{multicols}

    \flushright{
        {\phantom{qwe}}\\[0.5cm]
    }
    \centering{
        Санкт-Петербург\\
        2024
    }

    \vfill
\end{titlepage}

\Large
\tableofcontents
\newpage
\large

\section{Постановка задачи}

Интеграционное тестирование предполагает наличие нескольких модулей или, если приложение построено в соответствии с микросервисной архитектурой, микросервисов разработанного приложения. В качестве модулей могут быть функциональные части (например: модуль авторизации/аутентификации, модуль взаимодействия с пользователем, модуль интеграции со сторонним сервисом).

Необходимо выполнить интеграционное тестирование нескольких модулей в соответствии с предварительно описанными сценариями тестирования. Каждый сценарий тестирования должен представлять собой некоторый законченный вариант использования системы (кейс) пользователя.

Интеграционное тестирование предполагает командную работу над программным продуктом, поэтому необходимо проводить его с использованием одного из инструментов \texttt{CI}, доступных на рынке: \texttt{Gitlab}, \texttt{Jenkins}, \texttt{Bamboo}, \texttt{TeamCity}, etc.

Минимальное количество сценариев - 10.

\subsection{Обязательные требования}

\begin{enumerate}
    \item Предварительное формирование документа, описывающего тестовые сценарии.
    \item Применение заглушек (mock-сервисов или mock-объектов) для изоляции от окружения и внешних сервисов или для ускорения прохождения тестов.
    \item Рассмотрение негативных сценариев тестирования.
    \item Поднятие сервера непрерывной интеграции и запуск задачи по интеграционному тестированию по временному триггеру или по событию изменения кода приложения/интеграционных тестов.
\end{enumerate}

\subsection{Содержание отчета}

Отчет по интеграционному тестированию должен содержать:

\begin{enumerate}
    \item Отчет о выполненной работе, использованных инструментах, примененных техниках тест-дизайна.
    \item Тест-план со словесным описанием тестовых сценариев, которые планируется реализовать в интеграционном тестировании.
    \item Отчет о прохождении тестов с результатами на сервере непрерывной интеграции.
    \item Описание процедуры расширения тестового набора на примере добавления новой функциональной части (или модуля).
\end{enumerate}

\end{document}